% ============================================================
%  CHAPTER 1 — Introduction
%  Status : v1.0 — 2026-05-03
%  Logical chain: Problem → Gap → Proposed Solution
% ============================================================

\chapter{Introduction}
\label{chap:introduction}

% ─── 1.1 MOTIVATION ─────────────────────────────────────────
\section{Motivation}
\label{sec:motivation}

Humanoid robots occupy a singular position among robotic platforms:
their morphology — two legs, two arms, an upright torso — is shaped to
share the spaces that humans have built for themselves. Factories with
narrow aisles, hospitals with doorframes calibrated to wheelchair widths,
construction sites with uneven floors and overhead clearances: these
environments impose structural affordances that wheeled or
low-dimensional manipulators cannot exploit without fundamental redesign
of the surrounding infrastructure. A humanoid robot, in principle, can
walk through a corridor, open a door, and carry an object, relying on
the same geometry that human engineers assumed when the corridor was
laid out. This promise has motivated decades of research, from the
pioneering demonstrations of Honda's ASIMO and the HRP series to the
more recent platforms — Boston Dynamics Atlas, PAL Robotics TALOS,
Unitree H1 and G1 — which operate at human scale with torque-controlled
joints and on-board perception \cite{Koolen2016}.

Realising this promise, however, demands more than mechanical fidelity.
A 28--32 degree-of-freedom (DOF) humanoid body simultaneously executing
locomotion and a manipulation task faces a control problem of exceptional
difficulty: the system must coordinate leg joints to maintain dynamic
balance, arm and torso joints to position an end-effector, and the
shared trunk degrees of freedom that couple these two functions — all
at kilohertz control rates, subject to hard constraints on joint
torques, contact forces, and kinematic reachability
\cite{Sentis2005,Wensing2013}. Decomposing the robot into independent
subsystems — a leg controller and an arm controller — is the simplest
architectural choice, but it fails under precisely the conditions that
make humanoid platforms valuable: when locomotion perturbs the arm
workspace, when a heavy carried load shifts the whole-body centre of
mass, or when a manipulation task requires the robot to lean and thus
displace its support polygon \cite{Koolen2016}.

\WBC{} emerged as the theoretical answer to this coordination problem.
Rather than treating the limbs as independent, WBC considers all DOF as
a single dynamical system subject to a hierarchy of simultaneous tasks
and physical constraints \cite{Khatib1987,Sentis2005}. A postural
stability task maintains the robot upright; a locomotion task advances
the centre of mass; a manipulation task guides the end-effector to a
target pose; and each task contributes to the same joint torque command
through a priority-aware reconciliation mechanism. This unification
is what makes loco-manipulation — the simultaneous execution of
locomotion and manipulation — tractable in principle.

In practice, the gap between the theoretical framework and a deployable
implementation is substantial. Whole-body optimisation at joint-torque
level requires solving a constrained mathematical program at every
control step (typically 500\,Hz to 1\,kHz), incorporating contact
physics that are inherently uncertain, and doing so on processors that
must also run perception, state estimation, and safety monitoring. The
tension between constraint richness and computational tractability has
produced a stratified research landscape: lightweight task-priority
methods that sacrifice constraint handling for speed \cite{Sentis2005},
dense \QP{} formulations that capture contact physics rigorously but at
higher computational cost \cite{Wensing2013,Koolen2016}, and \MPC{}
layers that plan over a receding horizon but operate on reduced
dynamical models \cite{Orin2013}. Each layer solves a part of the
problem; no single, well-documented, reproducible framework has
integrated all three for simultaneous loco-manipulation and subjected
the integration to standardised evaluation.

This thesis is motivated by that gap.

\begin{criticalinsight}
  The argument from morphological fitness — humanoids for human spaces —
  carries rhetorical force but must not obscure the actual technical
  contribution. The motivation section establishes \emph{why the problem
  matters}, not why this particular solution is correct. A reviewer is
  entitled to ask whether the same task could be accomplished by a
  well-designed mobile manipulator; the honest answer is that, for a
  range of dexterous physical tasks, it could. The specific claim this
  thesis makes is narrower: that \emph{if} one has a humanoid robot, the
  absence of a reproducible, integrated WBC framework for simultaneous
  loco-manipulation is a scientific and engineering gap worth addressing
  rigorously. That claim is not weakened by the existence of alternative
  platforms.
\end{criticalinsight}


% ─── 1.2 PROBLEM STATEMENT ──────────────────────────────────
\section{Problem Statement}
\label{sec:problem}

The state of the art in \WBC{} for humanoid robots, as surveyed in
\Cref{chap:sota}, reveals two interacting deficiencies that the present
work targets directly.

\paragraph{The architectural coupling gap.}
A complete WBC pipeline for simultaneous loco-manipulation requires at
minimum two functional layers: a \emph{planning layer} that generates
kinematically and dynamically feasible reference trajectories (footstep
placements, centre-of-mass trajectories, end-effector targets), and an
\emph{execution layer} that translates these references into joint
torques while satisfying, at every control step, the constraints imposed
by contact physics, actuator limits, and whole-body dynamics.
\MPC{} operating on centroidal dynamics \cite{Orin2013} is the
established approach for the planning layer; \QP{}-based WBC is the
established approach for the execution layer \cite{Wensing2013,Koolen2016}.
Both are mature sub-fields with recognised algorithmic frameworks.
What remains underspecified, however, is the \emph{interface} between
these layers: how the contact schedule is passed from planner to
executor, how timing mismatches are handled when the MPC horizon and
the QP control loop run at different rates, and how constraint
violations at the execution level are fed back to the planner.
Existing implementations — Crocoddyl \cite{Crocoddyl}, TSID
\cite{DelPrete2016}, iCub-controllers \cite{Elobaid2023} — each resolve
this interface differently, often in platform-specific code that is not
documented at reproduction level. The field therefore lacks a consensus
architectural blueprint for the planning--execution coupling in
simultaneous loco-manipulation.

\paragraph{The evaluation standardisation gap.}
Even where WBC frameworks exist, their comparative assessment is
impeded by the absence of standardised benchmark scenarios and metrics.
Published results report different tasks (static balancing, walking,
push recovery, pick-and-place), different perturbation protocols
(impulse magnitude, direction, timing), and different success criteria —
often whether the robot ``fell'' is the only reported metric —
making it effectively impossible to determine whether a claimed
improvement is attributable to the controller design or to a more
favourable evaluation setup \cite{Sferrazza2024}. This evaluation
heterogeneity is particularly acute for loco-manipulation, where no
published benchmark addresses simultaneous locomotion and manipulation
with reproducible, quantified metrics.

This thesis addresses both gaps through two linked contributions:
(i)~an explicit design of the centroidal-\text{MPC}--\text{QP-WBC}
coupling, documented at reproduction level, and (ii)~a reproducible
evaluation protocol (scenarios S1 through S5) with standardised
metrics, validated against three baselines under identical simulation
conditions.

\begin{criticalinsight}
  The evaluation gap is arguably more consequential for the field's
  long-term progress than the architectural gap. A methodologically
  excellent framework that is evaluated idiosyncratically contributes
  little to cumulative science: subsequent groups cannot determine
  whether their own results represent progress. Conversely, a
  well-defined evaluation protocol, even paired with a modest
  implementation, creates a reusable reference point. The dual
  contribution structure of this thesis reflects this prioritisation:
  the framework is the vehicle; the reproducible protocol is the
  primary scientific deposit.
\end{criticalinsight}


% ─── 1.3 RESEARCH QUESTION ──────────────────────────────────
\section{Research Question and Objectives}
\label{sec:research-question}

The investigation is organised around the following primary research
question, locked as of 2026-04-30 in alignment with the thesis
supervisor:

\begin{quote}
  \textit{``How can a \WBC{} framework — combining a centroidal MPC
  planning layer with a \QP{}-based execution layer — be designed,
  implemented, and rigorously evaluated in simulation to enable a
  humanoid robot to perform stable, simultaneous locomotion and
  manipulation tasks under dynamic and contact constraints?''}
\end{quote}

\noindent Three sub-questions structure the investigation:

\begin{enumerate}[label=\textbf{SQ\arabic*}]
  \item What are the key limitations of existing WBC formulations
        (task-priority, QP-based, MPC-based) when applied to full-body
        humanoid coordination involving simultaneous locomotion and
        manipulation?
  \item How can contact dynamics be modelled and integrated in a
        simulation environment to faithfully replicate the physical
        behaviour expected of a torque-controlled humanoid?
  \item What scenarios and metrics constitute a rigorous, reproducible,
        simulation-based evaluation of a WBC framework for humanoid
        loco-manipulation?
\end{enumerate}

\noindent SQ1 is addressed in \Cref{chap:sota} through a critical
synthesis of the model-based WBC literature. SQ2 is addressed in
\Cref{chap:simulation} through the design of the MuJoCo simulation
environment and its contact model. SQ3 is addressed jointly in
\Cref{chap:methodology} (protocol design) and \Cref{chap:experiments}
(quantitative results).


% ─── 1.4 HYPOTHESES ─────────────────────────────────────────
\section{Hypotheses}
\label{sec:hypotheses}

The following three hypotheses operationalise the research question into
falsifiable empirical claims. Each hypothesis is associated with a
primary evaluation criterion; falsification conditions are stated
explicitly in \Cref{chap:methodology}.

\begin{hypothesis}{H1}
  A hierarchical \QP{}-based \WBC{} formulation, planned by a centroidal
  MPC layer, can achieve stable whole-body coordination across
  simultaneous locomotion and manipulation tasks — measured by a task
  success rate above a threshold defined from the benchmark literature
  and by bounded centre-of-mass tracking error throughout scenarios S1
  through S5.
\end{hypothesis}

\begin{hypothesis}{H2}
  MuJoCo simulation is a sufficient environment to validate
  proof-of-concept for the proposed WBC framework: tasks that succeed in
  simulation are indicative of hardware feasibility, modulo the
  sim-to-real gap characterised in \Cref{chap:discussion}.
\end{hypothesis}

\begin{hypothesis}{H3}
  Explicit contact modelling and task-priority hierarchies, as
  implemented in the proposed framework, produce lower centre-of-mass
  deviation and lower contact force variance than an unconstrained
  torque-control baseline under multi-contact transition scenarios.
\end{hypothesis}


% ─── 1.5 SCOPE AND LIMITATIONS ──────────────────────────────
\section{Scope and Limitations}
\label{sec:scope}

\begin{table}[h]
  \centering
  \caption{Thesis scope definition — in-scope and out-of-scope elements.}
  \label{tab:scope}
  \begin{tabular}{p{0.45\linewidth} p{0.45\linewidth}}
    \toprule
    \textbf{In Scope} & \textbf{Out of Scope} \\
    \midrule
    Simulation only (MuJoCo \cite{MuJoCo2012})
      & Hardware experiments \\
    WBC for simultaneous loco-manipulation
      & Manipulation-only or locomotion-only \\
    TALOS humanoid (PAL Robotics) URDF
      & Wheeled or legged non-humanoid robots \\
    Centroidal MPC + QP-based WBC
      & Deep RL as a primary contribution \\
    Reproducible evaluation protocol (S1--S5)
      & Sensor fusion / state estimation \\
    Comparative baselines: task-priority, torque ctrl, RL
      & Long-horizon task sequencing \\
    \bottomrule
  \end{tabular}
\end{table}

The scope boundary around simulation is a deliberate methodological
choice with two motivations. First, hardware experiments on a
full-scale torque-controlled humanoid require significant
infrastructure — a physical TALOS robot, custom actuator drivers,
safety systems, and operator expertise — that exceeds the resource
envelope of a solo master's thesis. Second, and more fundamentally,
simulation provides full state observability and exact reproducibility:
baselines can be evaluated on strictly identical conditions,
perturbations can be parameterised precisely, and results can be
generated from a public repository without hardware access. The value
of reproducibility is not a post-hoc rationalisation but a design
criterion: the evaluation gap identified in \Cref{sec:problem} is
partly a consequence of results that are hardware-dependent and
therefore non-reproducible by third parties.

The choice of MuJoCo \cite{MuJoCo2012} as the simulation environment is
motivated by its widely acknowledged contact stability, its adoption in
the model-based WBC community, and the availability of a validated
TALOS URDF model. The choice of TALOS as the target platform reflects
the same considerations: it is an open research platform with
documented kinematics and inertial parameters, torque-controlled joints
matched to the WBC execution layer's assumptions, and a growing body of
community software (Pinocchio \cite{Pinocchio}, Stack-of-Tasks).

The exclusion of deep reinforcement learning as a \emph{primary
contribution} does not imply a dismissal of RL-based approaches. Three
baselines include an RL policy, precisely to provide a cross-paradigm
comparison. The thesis does not argue that WBC is universally superior
to RL; it argues that the two approaches have different design
objectives — constraint satisfaction guarantees versus learned agility
— and that a fair comparison requires the standardised evaluation
protocol this work provides.


% ─── 1.6 CONTRIBUTIONS ──────────────────────────────────────
\section{Contributions}
\label{sec:contributions}

This thesis makes three principal contributions.

\paragraph{Contribution~1 — A documented WBC framework for simultaneous
loco-manipulation.}
A hierarchical control architecture is designed, implemented in
simulation, and documented at reproduction level. The architecture
couples a centroidal MPC planning layer — operating on a reduced dynamic
model of the TALOS humanoid — with a \QP{}-based whole-body execution
layer that enforces rigid-body dynamics, joint torque limits, and
friction-cone constraints at 1\,kHz. The interface between the two
layers, including contact schedule representation and rate decoupling,
is specified explicitly. The full implementation is released as an
open-access simulation repository (MuJoCo scene, TALOS URDF, controller
code, scenario configurations) to enable direct replication.

\paragraph{Contribution~2 — A reproducible evaluation protocol for
humanoid WBC.}
Five evaluation scenarios (S1:~static balance; S2:~bipedal walking;
S3:~push recovery; S4:~simultaneous walking and manipulation;
S5:~perturbation under loco-manipulation) are defined with precise
parameterisation of environment, initial conditions, perturbation
magnitude, and success criteria. A standardised metric set — CoM
tracking error, contact force variance, task success rate, and QP solve
time — is anchored in benchmark values from the literature. The
protocol is designed so that any WBC framework can be evaluated against
it given access to the simulation environment.

\paragraph{Contribution~3 — A comparative baseline analysis.}
The proposed framework is benchmarked against three baselines:
(i)~a task-priority null-space WBC without explicit inequality
constraints, (ii)~an unconstrained joint-torque controller, and
(iii)~a learned locomotion policy representative of the RL paradigm.
Results across S1--S5 provide a quantitative reference for future work
and allow hypotheses H1--H3 to be evaluated under controlled conditions.

The thesis does not claim to supersede established toolchains such as
Crocoddyl \cite{Crocoddyl} or TSID \cite{DelPrete2016}. The contribution
is complementary: an explicit coupling design for the specific case of
simultaneous loco-manipulation, and an evaluation infrastructure that
those toolchains could equally be assessed against.


% ─── 1.7 THESIS STRUCTURE ───────────────────────────────────
\section{Thesis Structure}
\label{sec:structure}

The remainder of this thesis is organised as follows.

\Cref{chap:sota} reviews the state of the art in \WBC{} for humanoid
robots. Model-based methods are examined in three groups — task-priority
and null-space formulations, QP-based WBC, and centroidal MPC — before
learning-based approaches are discussed as a baseline reference. The
chapter concludes with a synthesis that maps the literature gaps to the
contributions of this work.

\Cref{chap:methodology} formalises the proposed WBC framework. The
floating-base dynamics, task hierarchy, QP formulation, and centroidal
MPC layer are specified mathematically. The coupling interface and the
evaluation protocol (S1--S5) are defined in detail.

\Cref{chap:simulation} describes the simulation environment: the MuJoCo
scene, the TALOS URDF, the contact model, and the software stack
connecting Pinocchio's kinematic library \cite{Pinocchio} to the QP
solver.

\Cref{chap:experiments} presents quantitative results across scenarios
S1 through S5, with systematic comparison against the three baselines.

\Cref{chap:discussion} interprets the results, analyses the limits of
the simulation-based validation, and examines the sim-to-real gap in
the context of this framework.

\Cref{chap:conclusion} summarises the contributions, evaluates
hypotheses H1--H3 against the experimental evidence, and outlines
directions for future work, including hardware transfer and extensions
to long-horizon task sequencing.
